% Autoregressive analysis subsection; \input from 04_experiments.tex, before
% \subsection{Limitations}.
% Numbers from src/results/ar_baselines_tju.json (10 seeds/runs per cell),
% produced by src/ar_baselines_tju.py.  Delta values from
% src/diag_regen_vs_curvature.py.

\subsection{Autoregressive analysis}
\label{sec:ar}

Every result above is non-recursive: each prediction is made from a window of
true history. None of the three models compared here reports an autoregressive
evaluation either, since ours is the per-cycle protocol of
Table~\ref{tab:tableA}, PatchFormer's official inference is a full-sequence
non-recursive pass~\cite{ref_patchformer}, and RUL-Mamba's is likewise
\cite{ref_rulmamba}. Autoregression is therefore not a protocol on which any of
the three should be judged. It is, however, the question of what happens when
the model must consume its own output, which is what extrapolating past the last
observed cycle requires. We test it, and the outcome turns out not to be a
property of the architectures at all.

\paragraph{Protocol.}
We launch at $\mathrm{EOL}-k$ for $k\in\{50,100\}$, warm up on the $W$ true
cycles preceding the launch, and thereafter feed the model only its own
predictions. A fixed horizon is what makes the comparison meaningful: launching
at SP instead would give CALCE a 340-cycle horizon and MIT a 289-cycle one,
confounding the dataset with the horizon length. We report the crossing-cycle
error AE together with its ratio to the horizon, the trajectory MAE and MSE over
the interval from the launch to the \emph{true} EOL, and the predicted capacity
at the true EOL. We deliberately do not report a crossing rate. On CALCE our own
model crosses in 7 of 10 runs while missing the crossing cycle by twice the
horizon, so a crossing count would score that as a success.

\begin{table}[htbp]
\centering
\caption{Autoregressive rollout, ten seeds/runs per cell, each model at its
smallest SP and the dataset's native window. AE is the median crossing-cycle
error over the runs that crossed, and `--' marks a model that never crossed.
MAE and MSE are trajectory errors in normalized capacity, taken over the
launch-to-true-EOL interval. ``EOL err'' is the absolute error of the predicted
capacity at the true EOL, as a percentage of rated capacity.}
\label{tab:ar}
\scriptsize
\setlength{\tabcolsep}{2pt}
\begin{tabular}{lll r rr r r r}
\toprule
Dataset & Launch & Model & Crossed & AE & AE/$k$ & MAE & MSE & EOL err \\
\midrule
CALCE   & EOL-50  & Ours & 7/10 & 102 & 2.04 & 0.0165 & 6.6e-4 & 7.0\% \\
        &         & PF   & 1/10 &  56 & 1.12 & 0.0187 & 9.1e-4 & 7.7\% \\
        &         & RM   & 0/10 &  -- &   -- & 0.0199 & 9.6e-4 & 8.1\% \\
        & EOL-100 & Ours & 7/10 &  55 & 0.55 & 0.0189 & 6.6e-4 & 6.1\% \\
        &         & PF   & 1/10 &  33 & 0.33 & 0.0223 & 1.1e-3 & 9.4\% \\
        &         & RM   & 0/10 &  -- &   -- & 0.0267 & 1.4e-3 & 10.8\% \\
\midrule
MIT     & EOL-50  & Ours & 10/10 &  2 & 0.04 & 0.0039 & 5.3e-5 & 0.74\% \\
        &         & PF    &  8/10 & 12 & 0.25 & 0.0168 & 8.1e-4 & 4.9\% \\
        &         & RM    & 10/10 &  2 & 0.05 & 0.0044 & 2.4e-5 & 0.48\% \\
        & EOL-100 & Ours &  7/10 & 12 & 0.12 & 0.0135 & 4.1e-4 & 3.4\% \\
        &         & PF    &  7/10 & 24 & 0.24 & 0.0360 & 5.3e-3 & 12.5\% \\
        &         & RM    &  0/10 & -- &   -- & 0.0176 & 4.9e-4 & 4.1\% \\
\bottomrule
\end{tabular}
\end{table}

On MIT all three models roll out (Table~\ref{tab:ar}). Over a 50-cycle horizon
our model and RUL-Mamba both predict the EOL capacity to within 0.8\% of rated
with a median crossing error of two cycles, and PatchFormer, whose trajectory
MAE is four times larger, still crosses in 8 of 10 runs. At 100 cycles
RUL-Mamba freezes in all ten runs while the other two hold.

On CALCE none of them works. Every model lands 6.1--10.8\% of rated capacity
away from the true EOL value, and 7\% is not a usable forecast. The failure
\emph{modes} differ. Ours is bimodal, with four seeds freezing and six diverging
and nothing in between, while PatchFormer and RUL-Mamba freeze in all ten runs
each. The outcome does not.

\paragraph{Why the two datasets separate.}
The decoder here is window-relative: the network emits a scalar $z_t$ and the
prediction is $\hat y_t = z_t\,\sigma(w_t) + \mu(w_t)$, the window mean pushed
out by $z$ times the window's spread. For a window decaying at rate $r$, the
value a perfect one-step model must emit is
$z_{\mathrm{req}} = -\tfrac{W+1}{2}\,(r/\sigma)$, while reproducing a steady
decay requires the constant
$z_{\mathrm{crit}} = -\tfrac{W+1}{2}\,A^{-1}$ with $A=\sqrt{W(W+1)/12}$, the
unbiased standard deviation of a ramp and therefore the estimator the decoder
itself uses. The rate cancels in the latter, so the fixed point is set by the
decoder alone and not by the battery. Their difference
\begin{equation}
\delta \;=\; z_{\mathrm{req}} - z_{\mathrm{crit}}
\;=\; \tfrac{W+1}{2}\left(\tfrac{1}{A} - \tfrac{r}{\sigma}\right) \;\ge\; 0
\label{eq:delta}
\end{equation}
is non-negative for any noisy window and measures how far the decode sits from
its fixed point. A constant departure above that point makes the rollout's decay
rate fall geometrically, with half-life
$n_{1/2} \approx 0.693\,W/(\sqrt{3}\,\delta)$. Both terms come from the series
alone, before any model is trained. At $W{=}64$,
$\delta(\mathrm{MIT},\mathrm{EOL}{-}50) = +0.016$, a half-life of 1630 cycles
against a 50-cycle horizon, while
$\delta(\mathrm{CALCE},\mathrm{EOL}{-}50) = +0.238$, a half-life of 108 cycles,
inside the horizon. The criterion separates the two datasets, and the separation
is independent of which of the three models is rolled out.

\paragraph{Scope.}
The derivation holds $z$ fixed, whereas the model's $z$ is a function of the
window; $\delta$ is therefore a local quantity read at the launch point, not a
closed-form prediction of a particular checkpoint. It governs freeze and
divergence, not the accumulation of one-step error that dominates as the horizon
grows, which is what the EOL-100 rows of Table~\ref{tab:ar} show relative to the
EOL-50 rows. It should be read as a screening test for the failure modes above,
not as a guarantee of accuracy. Two cautions on its magnitude. The two terms in
Eq.~\eqref{eq:delta} agree to within a few tenths of a percent on a well-behaved
cell, so $\delta$ is numerically fragile and should be read for its sign and its
ordering rather than its value. And where they nearly cancel entirely, as on
NASA at $W{=}30$, $\delta$ varies with the launch point by as much as its own
range and predicts nothing. The separation above is demonstrated on two datasets
whose $\delta$ differs by an order of magnitude; we state it as a consistent
screening criterion, not as a validated law.

\paragraph{Implication.}
Autoregressive rollout should not be used as an evaluation protocol for these
models: whether it works is decided by the window-relative decode and by the
data, and all three architectures succeed or fail together. The value of
$\delta$ is as a screening quantity, computable from the series before training;
it identifies the datasets on which rollout is worth attempting at all.
